# Chapter Extension: Ge-on-Si PIN Photodetector
## Based on: Pilori (2019), *DSP Techniques for High-Speed Optical Communications Links*, Politecnico di Torino

---

## PART I — EXTRACTED INSIGHTS (Condensed, Technical)

### A. Physical Modelling

**Square-law detection and photocurrent linearity.** The attached thesis (Ch. 4) makes explicit a fundamental but frequently elided point: direct detection is a *square-law* process. The photocurrent is proportional to the squared modulus of the optical field envelope, not to the field amplitude itself. For a field carrying a complex information-bearing envelope $s(t)$ superposed on a carrier $c$, the photocurrent after an ideal photodiode takes the form $i(t) = |s(t) + c|^2$, which expands into three distinct terms: the signal power $|s(t)|^2$, a DC bias $c^2$, and a cross-beat $2c \cdot \mathrm{Re}\{s(t)\}$ that constitutes the useful, linearly demodulated signal. The linear approximation $I_\mathrm{ph} = \mathcal{R} P_\mathrm{opt}$—which the existing chapter applies throughout its metric definitions—is valid precisely because the carrier power dominates: $c^2 \gg |s(t)|^2$, so that $|s(t)|^2$ (the signal–signal beating interference, SSBI) is negligible. This approximation fails for high-modulation-depth signals and for the analysis of photodetector non-linearity, and its domain of validity must be stated explicitly.

**Chromatic dispersion and the direct-detection transfer function.** When dispersion is present, the electric field acquires a frequency-dependent phase $H_\mathrm{CD}(f) = \exp(-j2\pi^2\beta_2 f^2 z)$ before square-law detection. Substituting into the square-law relation generates a baseband electrical transfer function $\cos(\pi\beta_2 f^2 z)$ with periodic deep notches at frequencies $f_\mathrm{notch} = \sqrt{k/(|\beta_2|z)}$, $k = 1, 2, \ldots$. These notches destroy spectral components of the information signal irrecoverably at the output of the photodetector. This analysis is the precise quantitative reason why the O-band ($\lambda \approx 1310$ nm, where $\beta_2 \approx 0$ for standard SMF) must be used for dispersion-uncompensated direct-detection links: the notch frequencies are pushed to infinity, eliminating the transfer-function penalty entirely.

**Noise current spectral density as an engineering metric.** The thermal noise of a practical transimpedance-amplifier (TIA)–photodiode assembly is most compactly expressed through the input-referred equivalent noise current density $i_n$ [A$/\sqrt{\mathrm{Hz}}$], which encapsulates both the thermal noise of the load and the amplifier's internal noise. The total noise variance integrated over detection bandwidth $B_e$ is $\sigma_n^2 = i_n^2 B_e$ [A$^2$], a form directly readable from component datasheets and directly comparable with the shot-noise contribution $2q(I_d + I_\mathrm{ph}) B_e$.

**PAM-$M$ BER and the OMA framework.** For an $M$-level pulse-amplitude-modulated direct-detection link dominated by Gaussian noise of variance $\sigma_n^2$, the symbol-error rate is determined entirely by the ratio of the electrical eye opening to the noise standard deviation. The eye opening is controlled not by average optical power but by the outer optical modulation amplitude $\mathrm{OMA}_\mathrm{out} = P_{M-1} - P_0$, which is independent of the extinction ratio and is the Ethernet-standard figure of merit. The BER framework expressed in terms of $\mathrm{OMA}_\mathrm{out}$, $\mathcal{R}$, and $i_n$ is strictly more general than the OOK shot-noise formula of the existing chapter because it captures the $\sim 3$ dB/bit OMA penalty incurred when the modulation order is doubled, and it explicitly separates the optical design variable ($\mathrm{OMA}_\mathrm{out}$) from the electrical noise variable ($i_n$).

**Photodetector bandwidth limitation as ISI.** When the photodetector's $f_{3\,\mathrm{dB}}$ falls below the Nyquist frequency of the transmitted symbol stream, the device's low-pass transfer function introduces inter-symbol interference (ISI) at the decision point. For a symbol rate $R_s$, the Nyquist frequency is $f_\mathrm{Nyq} = R_s/2$, and the bandwidth utilisation efficiency condition $f_\mathrm{Nyq} \leq 0.7\,f_{3\,\mathrm{dB}}$ (the single-pole passband approximation) must be satisfied for error-free detection without equalisation. Beyond this regime, a feed-forward equaliser (FFE) or decision-feedback equaliser (DFE) must restore the roll-off. The FFE operates by convolving the received sequence with a finite-impulse-response filter that approximately inverts the photodetector's low-pass response; its complexity (tap count) scales with the ratio $f_\mathrm{Nyq}/f_{3\,\mathrm{dB}}$, so that $f_{3\,\mathrm{dB}}$ directly controls post-detector DSP cost.

### B. Numerical Methods

**Transfer-function factorisation as a simulation shortcut.** The Pilori thesis treats the photodetector's frequency response as a cascade of independently computed LTI stages: transit-time sinc roll-off and RC pole. This factorisation is adopted here as the analytical bandwidth model of the existing chapter's Eq. (bw\_total), but can be extended to include the third-order roll-off of realistic bondpad parasitics, the dispersive phase factor from chromatic dispersion, and the finite TIA bandwidth. Each stage is independently measured or extracted, and the cascade product gives the end-to-end link transfer function.

**Equivalent SNR for multi-carrier modulation.** For any multi-carrier format (OFDM, DMT, SCM) applied to a direct-detection channel, the achievable information rate is determined by the geometric mean of per-subcarrier SNRs: $\mathrm{SNR}_\mathrm{eq} = \left(\prod_{i=1}^N \mathrm{SNR}_i\right)^{1/N}$. The geometric mean arises because each subcarrier independently satisfies the Shannon bound, so the composite rate is the sum of per-subcarrier rates. This formulation generalises the single-carrier sensitivity of the existing chapter to wideband modulation formats and provides a rigorous framework for evaluating the impact of the photodetector's frequency-selective roll-off on aggregate capacity.

### C. Device-Level Insights

**Condition for linear photocurrent approximation.** The linear relation $I_\mathrm{ph} = \mathcal{R} P_\mathrm{opt}$ holds only when the SSBI term $|s(t)|^2$ is small relative to the cross-beat $2c \cdot \mathrm{Re}\{s(t)\}$, i.e., when $\mathrm{CSPR} \equiv c^2/\sigma_s^2 \gg 1$. In a CW-illuminated photodetector with modulated power $P(t) = |s(t)|^2$, the approximation is equivalent to a large extinction ratio, and the small-signal responsivity formula holds. For high-modulation-depth or multi-level formats, higher-order corrections introduce non-linearity in the photocurrent characteristic that degrades the IQE measurement and must be accounted for in the responsivity extraction.

**O-band vs. C-band for direct detection.** The zero-dispersion wavelength of standard SMF lies within the O-band ($\approx 1310$ nm), yielding $|\beta_2| < 2$ ps$^2$/km. At this dispersion level, the first transfer-function notch for a 10 Gbit/s OOK signal occurs beyond 100 km, completely outside the intra-data-centre reach of interest. Silicon's transparency at 1310 nm, which necessitates the Ge absorber on physical grounds (established in the existing chapter through the band-structure argument), therefore coincides precisely with the wavelength choice that maximises direct-detection system margin. The two motivations—material absorption physics and dispersion-free transmission—are thus mutually reinforcing constraints, not independent decisions.

### D. Representation Techniques

The Pilori thesis employs the following representation strategies that can be adopted as extensions:

- BER curves as a function of outer OMA (in dBm) rather than received optical power, which separates modulator linearity from detector noise.
- Equivalent SNR vs. CSPR curves for the quantification of shot-noise–SSBI trade-off.
- Transfer-function magnitude vs. frequency (in GHz) showing independent transit-time, RC, and combined roll-offs, with the $-3$ dB point and Nyquist frequency annotated.
- Required OSNR vs. dispersion-link distance for DSB/SSB modulation comparison.

---

## PART II — MAPPING TO CHAPTER SECTIONS

| Extension ID | Insertion Point | Content Added |
|---|---|---|
| **E1** | Section 1.1, after Eq. (beer\_lambert) | Square-law detection: formal derivation of photocurrent from Poynting vector; domain of validity of linear approximation |
| **E2** | Section 1.4.3, after Eq. (bw\_rc) | Photodetector as LTI channel; ISI onset condition; equaliser tap-count scaling |
| **E3** | Section 1.5, between Eq. (thermal) and Eq. (nep) | Input-referred noise current density; unified TIA–photodiode noise model; total noise variance |
| **E4** | Section 4.2, after the PAM-4 sentence | PAM-$M$ SER/BER framework via OMA; per-bit OMA penalty; chromatic dispersion transfer function; O-band rationale |
| **E5** | Section 3.2, after Eq. (fdtd\_abs) | FDTD E-field to photocurrent: connection between Poynting intensity monitor and square-law detection; regime of linear operation |

---

## PART III — EXTENDED PARAGRAPHS

---

### [Extended Paragraph — Section 1.1, after Eq. (beer\_lambert)]

**Insert after the sentence ending "…rather than an independent empirical postulate." and the absorbed power density equation (abs\_density), before the figure and before Section 1.2.**

A corollary of the plane-wave analysis is the identification of the photodetection mechanism as a fundamentally square-law operation on the optical field. The quantity $I(z) = \tfrac{1}{2}\mathrm{Re}[\tilde{\mathbf{E}} \times \tilde{\mathbf{H}}^*]\cdot\hat{z}$, to which the Beer-Lambert law applies, is proportional to $|\tilde{\mathbf{E}}|^2$, not to $|\tilde{\mathbf{E}}|$ itself. A photodiode converts this time-averaged optical intensity into a terminal current by absorbing photons and generating electron-hole pairs at the rate of \cref{eq:gen_rate}; because the intensity is quadratic in the field amplitude, the resulting photocurrent $i_\mathrm{ph}(t)$ is inherently a square-law functional of the optical field envelope. Formally, for a quasi-monochromatic field at carrier frequency $\omega_0$ with a slowly varying complex envelope $\tilde{s}(t)$,
\begin{equation}
    \tilde{\mathbf{E}}(t) = \tfrac{1}{2}\tilde{s}(t)\,e^{-i\omega_0 t}\hat{\mathbf{e}} + \mathrm{c.c.},
    \label{eq:envelope_field}
\end{equation}
the time-averaged Poynting intensity at the detector facet is $I(t) = |\tilde{s}(t)|^2/(2\eta_0)$, where $\eta_0 = \sqrt{\mu_0/\varepsilon_0}$ is the impedance of free space. The responsivity $\mathcal{R}$ of \cref{eq:responsivity} then defines the photocurrent strictly as
\begin{equation}
    i_\mathrm{ph}(t) = \mathcal{R}\cdot\frac{|\tilde{s}(t)|^2}{2\eta_0} = \mathcal{R}\cdot P(t),
    \label{eq:square_law}
\end{equation}
where $P(t) = |\tilde{s}(t)|^2/(2\eta_0)$ is the instantaneous optical power. The linearity of \cref{eq:square_law} is therefore a linearity in optical power, not in optical field amplitude. For a continuous-wave signal of constant power $P_\mathrm{opt}$, \cref{eq:square_law} reduces identically to the responsivity definition of \cref{eq:responsivity}. For an intensity-modulated signal, however, a non-trivial consequence of the square-law relation emerges when the field contains both a data-bearing component $\tilde{u}(t)$ and an unmodulated component $c$ (e.g., a residual DC bias or carrier): the total field envelope is $\tilde{s}(t) = \tilde{u}(t) + c$, and the photocurrent becomes
\begin{equation}
    i(t) = \mathcal{R}\bigl[|\tilde{u}(t)|^2 + c^2 + 2c\,\mathrm{Re}\{\tilde{u}(t)\}\bigr].
    \label{eq:square_law_expand}
\end{equation}
The term $2c\,\mathrm{Re}\{\tilde{u}(t)\}$ constitutes the linearly demodulated signal; the term $|\tilde{u}(t)|^2$ is the signal-signal beating interference (SSBI), which is a baseband distortion product of the square-law process; and $c^2$ is a DC offset. The linear photocurrent approximation $i_\mathrm{ph} \approx \mathcal{R}\cdot P_\mathrm{opt}$, employed throughout the metric framework of \cref{subsec:pd_metrics}, is valid precisely when $c^2 \gg |\tilde{u}(t)|^2$—equivalently, when the carrier-to-signal power ratio $\mathrm{CSPR} \equiv c^2/\langle|\tilde{u}|^2\rangle \gg 1$—so that SSBI is negligible relative to the cross-beat term. In the Ge-on-Si waveguide photodetector under CW illumination at $P_\mathrm{opt} = \SI{100}{\micro\watt}$, this condition is satisfied by construction, and \cref{eq:responsivity} is an exact definition in the linear regime. The square-law nature of \cref{eq:square_law} becomes quantitatively consequential for high-modulation-depth multilevel formats (PAM-4, PAM-8) and determines the onset of non-linear photocurrent distortion that must be incorporated into any system-level link budget. The responsivity and EQE definitions of \cref{subsec:pd_metrics} remain valid in the linear regime; extensions to the non-linear regime require higher-order Volterra-series modelling that is beyond the scope of this chapter.

---

### [Extended Paragraph — Section 1.4.3 (RC Limit), after Eq. (bw\_rc)]

**Insert after the paragraph ending "…is the principal motivation for the U-shaped electrode architecture discussed in \cref{subsec:pd_architecture}." and before Section 1.4.4 (Capacitance and Series Resistance).**

The combined transit-time and RC bandwidth model of \cref{eq:bw_total} admits a compact linear-systems representation that makes the impact on communication performance immediately transparent. Define the small-signal photocurrent transfer function $H_\mathrm{PD}(f)$ as the ratio of the photocurrent amplitude at frequency $f$ to its DC value, normalised to unity at $f = 0$. Under the factorisation assumption that transit-time and RC roll-offs are independent, $H_\mathrm{PD}(f)$ is well-approximated by the product
\begin{equation}
    H_\mathrm{PD}(f)
    = \underbrace{\mathrm{sinc}\!\left(\frac{f}{f_\mathrm{tt}}\right)}_{H_\mathrm{tt}(f)}
    \cdot\underbrace{\frac{1}{1+i\,f/f_\mathrm{RC}}}_{H_\mathrm{RC}(f)},
    \label{eq:H_pd}
\end{equation}
where $f_\mathrm{tt}$ and $f_\mathrm{RC}$ are defined by \cref{eq:bw_tt,eq:bw_rc} respectively, $\mathrm{sinc}(x) = \sin(\pi x)/(\pi x)$, and the combined $-3\,\mathrm{dB}$ bandwidth $f_{3\,\mathrm{dB}}$ of \cref{eq:bw_total} is the frequency at which $|H_\mathrm{PD}(f)|^2 = 1/2$. The Fourier magnitude $|H_\mathrm{tt}(f)|$ is the normalised spectrum of the rectangular current pulse of duration $\tau_\mathrm{tr} = W_i/v_\mathrm{d,sat,h}$; the first zero of $|H_\mathrm{tt}(f)|$ occurs at $f = f_\mathrm{tt}/0.45 \approx \SI{113}{\giga\hertz}$ for the present device, well above the operating bandwidth. The single-pole $H_\mathrm{RC}(f)$ dominates the roll-off beyond $f_\mathrm{RC}$, and the combined $|H_\mathrm{PD}(f)|^2$ determines the effective bandwidth of the photocurrent channel that the subsequent receiver electronics must process. A transmitted symbol stream at rate $R_s$ has a Nyquist bandwidth $f_\mathrm{Nyq} = R_s/2$. The condition for ISI-free reception without equalisation is $f_\mathrm{Nyq} \leq \beta\,f_{3\,\mathrm{dB}}$, where the dimensionless factor $\beta \approx 0.7$ accounts for the spectral roll-off of the single-pole model, ensuring that the signal power at the Nyquist frequency has degraded by no more than $\sim\SI{3}{\decibel}$ relative to the DC responsivity. Violation of this condition—equivalently, $R_s > 2\beta\,f_{3\,\mathrm{dB}}$—introduces ISI whose mitigation requires a feed-forward equaliser (FFE) at the receiver. The FFE tap count required to compensate a roll-off deficit of $\Delta f = f_\mathrm{Nyq} - \beta\,f_{3\,\mathrm{dB}}$ scales approximately as $N_\mathrm{tap} \sim 2f_\mathrm{Nyq}/f_{3\,\mathrm{dB}}$, so that every doubling of $f_{3\,\mathrm{dB}}$ relative to $f_\mathrm{Nyq}$ halves the DSP complexity. The present device, with $f_{3\,\mathrm{dB}} \geq \SI{103}{\giga\hertz}$ and the intended operating rate of $R_s = \SI{100}{\giga\baud}$ (giving $f_\mathrm{Nyq} = \SI{50}{\giga\hertz}$), satisfies $f_\mathrm{Nyq}/f_{3\,\mathrm{dB}} \approx 0.49 < \beta$, placing the system firmly within the ISI-free regime and eliminating any equaliser requirement directly attributable to photodetector bandwidth limitation. The transfer function $H_\mathrm{PD}(f)$ also serves as the device-level input to the link equaliser model developed in the digital signal processing chapters of this thesis, where it enters as the analogue channel model that the equaliser inversion filter must reproduce.

---

### [Extended Paragraph — Section 1.5 (Noise Spectral Densities and NEP), between Eq. (thermal) and Eq. (nep)]

**Insert after the Johnson noise equation (thermal) and before the NEP definition (nep).**

In practice, the photodetector is always followed by a transimpedance amplifier (TIA), whose own internal noise—amplifier input-referred thermal noise, channel noise, and gate leakage in the case of MOSFET-input stages—must be added to the detector-intrinsic noise terms of \cref{eq:shot,eq:thermal}. The combined noise of the photodetector–TIA assembly is most compactly characterised through the one-sided input-referred equivalent noise current spectral density:
\begin{equation}
    i_n^2 = S_\mathrm{shot} + S_\mathrm{thermal} + S_\mathrm{TIA}
    \quad [\si{\ampere\squared\per\hertz}],
    \label{eq:in_total}
\end{equation}
where $S_\mathrm{TIA}$ is the TIA's input-referred noise current PSD, which for a resistive-feedback transimpedance stage is $S_\mathrm{TIA} = 4kT/R_F + (2\pi f)^2 C_\mathrm{in}^2 S_{v,\mathrm{amp}}$, with $R_F$ the feedback resistance, $C_\mathrm{in} = C_j + C_\mathrm{par} + C_\mathrm{gate}$ the total input capacitance, and $S_{v,\mathrm{amp}}$ the amplifier's input voltage noise PSD. The quantity $i_n = \sqrt{i_n^2}$ [A$/\sqrt{\mathrm{Hz}}$] is directly reported on component datasheets and represents the minimum detectable current per unit electrical bandwidth. Integrating \cref{eq:in_total} over the detection bandwidth $B_e = f_{3\,\mathrm{dB}}$ gives the total noise variance at the receiver input:
\begin{equation}
    \sigma_n^2 = i_n^2\,B_e
    \quad [\si{\ampere\squared}],
    \label{eq:sigma_n}
\end{equation}
a compact expression that consolidates all noise contributions into a single measurable scalar. The shot-noise-only limit, obtained by setting $S_\mathrm{thermal} = S_\mathrm{TIA} = 0$ and $I_\mathrm{ph} \ll I_d$, gives $i_n = \sqrt{2qI_d}$, which for the simulated $I_d = \SI{1.3}{\nano\ampere}$ evaluates to $i_n \approx \SI{0.645}{\pico\ampere\per\sqrt\hertz}$. Practical TIA stages realised in \SI{90}{\nano\metre} CMOS technology achieve input-referred noise densities in the range $\SI{10}{\pico\ampere\per\sqrt\hertz}$ to $\SI{20}{\pico\ampere\per\sqrt\hertz}$ over bandwidths exceeding \SI{100}{\giga\hertz}, so that $S_\mathrm{TIA}$ dominates $S_\mathrm{shot}$ at the dark-current level of the present device by approximately two orders of magnitude. The dominance of TIA noise over photodetector shot noise in this configuration implies that further reduction of $I_d$ below \SI{1.3}{\nano\ampere} yields negligible improvement in the total $\sigma_n^2$, and that the receiver sensitivity is TIA-limited rather than detector-limited—a conclusion with direct implications for the design prioritisation in the transceiver chain described in the system chapters. The NEP of \cref{eq:nep} and the specific detectivity of \cref{eq:dstar} are intrinsic device metrics that quantify the minimum detectable power in the absence of external electronics; the more operationally relevant figure of merit for the integrated receiver is the optical sensitivity $P_\mathrm{sens}$ of \cref{sec:pd_system}, which is derived from the total $\sigma_n$ of \cref{eq:sigma_n} after substituting the assembly noise current.

---

### [Extended Paragraph — Section 4.2 (Bandwidth and Modulation Format Implications), after the PAM-4 sentence]

**Insert after the sentence ending "…the same $f_{3\,\mathrm{dB}}$ supports up to approximately \SI{200}{\giga\bit\per\second} with reduced bandwidth penalty." and before the DSP-layer equalisation sentence.**

A more precise quantification of the modulation-order trade-off is obtained through the PAM-$M$ symbol-error-rate framework, which generalises the OOK sensitivity expression of \cref{subsec:pd_sens} to multilevel intensity modulation. For an $M$-level PAM signal transmitted over a direct-detection channel with additive Gaussian noise of variance $\sigma_n^2 = i_n^2 B_e$ (following the total noise model of \cref{eq:in_total,eq:sigma_n}), and detected by a photodiode of responsivity $\mathcal{R}$, the symbol-error ratio is determined by the ratio of the inter-level spacing to the noise standard deviation. Defining the outer optical modulation amplitude $\mathrm{OMA}_\mathrm{out} = P_{M-1} - P_0$ as the difference between the uppermost and lowermost transmitted power levels—a quantity independent of the extinction ratio and directly controlled by the modulator drive amplitude—the inter-level spacing is $d = \mathcal{R}\,\mathrm{OMA}_\mathrm{out}/(M-1)$, and the symbol-error rate is:
\begin{equation}
    P_s = \frac{M-1}{M}
    \,\mathrm{erfc}\!\left(
    \frac{\mathcal{R}\cdot\mathrm{OMA}_\mathrm{out}}
         {\sqrt{8\,B_e\,(M-1)}\,i_n}
    \right).
    \label{eq:ser_pam}
\end{equation}
The argument of the erfc function—the detection SNR—is directly proportional to $\mathcal{R}\cdot\mathrm{OMA}_\mathrm{out}/i_n$, confirming that responsivity, modulator swing, and receiver noise enter the BER on equal footing. Holding $P_s$ constant at a target BER and doubling the modulation order from $M$ to $2M$ requires an approximately $\SI{3}{\decibel}$ increase in $\mathrm{OMA}_\mathrm{out}$ to maintain the same argument of the erfc; this is the per-bit OMA penalty that accumulates linearly in dB as the constellation order is increased. For the present device with $\mathcal{R} = \SI{0.95}{\ampere\per\watt}$ and the TIA-dominated noise floor of $\sim\SI{15}{\pico\ampere\per\sqrt\hertz}$ at the full $B_e = \SI{103}{\giga\hertz}$, the required OMA for PAM-2 at BER $= 10^{-4}$ is approximately $\SI{-18}{\dBm}$, while PAM-4 at half the symbol rate (and hence $B_e/2 = \SI{51.5}{\giga\hertz}$, which halves the noise variance) requires approximately $\SI{-12}{\dBm}$—a net sensitivity penalty of $\SI{6}{\decibel}$ despite the $\SI{3}{\decibel}$ noise reduction from the narrower electrical bandwidth. The additional degradation beyond the simple $\SI{3}{\decibel}$/bit law arises from the tighter eye opening of the 4-level constellation relative to the 2-level case at fixed $\mathrm{OMA}_\mathrm{out}$. A second consideration that constrains the usable bandwidth of the photodetector in the context of a dispersive channel concerns the small-signal transfer function of direct detection in the presence of group-velocity dispersion. Chromatic dispersion in a single-mode fibre is modelled as an all-pass optical filter with transfer function $H_\mathrm{CD}(f) = \exp(-j\pi\beta_2 f^2 z)$, where $f$ is the baseband frequency, $z$ the propagation distance, and $\beta_2$ [ps$^2$/km] the group-velocity dispersion parameter. After propagation through this filter and square-law detection, the small-signal electrical transfer function takes the form $\mathrm{Re}\{H_\mathrm{CD}(f)\} = \cos(\pi\beta_2 f^2 z)$, which exhibits deep nulls at frequencies satisfying $\pi|\beta_2|f^2 z = \pi/2 + k\pi$, i.e., $f_k = \sqrt{(2k+1)/(2|\beta_2|z)}$ for $k = 0, 1, 2, \ldots$. These nulls completely destroy spectral components of the photocurrent at those frequencies—a loss that no subsequent equalisation can recover because the information is irreversibly annihilated at the detector output. For standard SMF near the O-band zero-dispersion wavelength ($|\beta_2| \approx 0.5$ ps$^2$/km at \SI{1310}{\nano\metre}), the first null at a propagation distance of $z = \SI{1}{\kilo\metre}$ occurs at $f_0 = \sqrt{1/(2\times 0.5\times 10^3)} \approx \SI{1}{\tera\hertz}$, far outside any conceivable electrical bandwidth. The photodetector at \SI{1310}{\nano\metre} therefore operates in a dispersion-null-free regime for all intra-datacenter reach distances, a condition that, combined with the band-structure argument of \cref{subsec:pd_bandstructure}, makes O-band operation the uniquely correct choice for both material-physics and systems-level reasons simultaneously.

---

### [Extended Paragraph — Section 3.2 (Optical Simulation: FDTD), after Eq. (fdtd\_abs)]

**Insert after the mesh-convergence sentence ending "…less than \SI{0.5}{\percent} deviation in $\mathcal{A}$." and before the absorber length sweep sentence.**

A physically essential clarification connecting the FDTD electromagnetic solution to the photocurrent calculation concerns the precise quantity that the power monitors extract. The time-domain FDTD solver advances the real-valued field components $\mathbf{E}(\mathbf{r},t)$ and $\mathbf{H}(\mathbf{r},t)$ on the Yee staggered grid, and the frequency-domain power recorded by the monitors is the time-averaged Poynting vector $\mathbf{S}(\mathbf{r},\omega) = \tfrac{1}{2}\mathrm{Re}[\tilde{\mathbf{E}} \times \tilde{\mathbf{H}}^*]$, which is quadratic in the field amplitudes $\tilde{\mathbf{E}}$ and $\tilde{\mathbf{H}}$. The absorbed power density of \cref{eq:abs_density}, computed as $\mathcal{P}_\mathrm{abs} = -\nabla\cdot\mathbf{S}$, is therefore also quadratic in the field—consistent with the square-law photodetection mechanism formalised in \cref{eq:square_law}. The optical generation rate $G_\mathrm{opt}(\mathbf{r}) = \mathcal{P}_\mathrm{abs}(\mathbf{r})/\hbar\omega$ of \cref{eq:gen_rate}, transferred to the CHARGE solver via the \texttt{ge\_gen\_oband.mat} file, is thus already expressed in the physically correct power-quadratic form that maps directly to the photocurrent via $I_\mathrm{ph} = q\int_\Omega G_\mathrm{opt}\,\mathrm{d}V$. The linear photocurrent approximation $I_\mathrm{ph} = \mathcal{R}\,P_\mathrm{opt}$ is implicit in this computation: because the FDTD simulation is performed at a fixed, low incident power $P_\mathrm{opt} = \SI{100}{\micro\watt}$, the optical field within the germanium absorber remains well within the unsaturated, linear response regime, where carrier densities $\Delta n$, $\Delta p$ injected by illumination are negligibly small compared to the doping-controlled equilibrium concentrations. This can be verified explicitly: the total carrier generation rate $\int_\Omega G_\mathrm{opt}\,\mathrm{d}V \approx I_\mathrm{ph}/q \approx \SI{5.9e11}{\per\second}$, and at steady state the photo-injected excess carrier density in the intrinsic region is $\Delta n \approx G_\mathrm{opt,avg}\,\tau_0 \sim \SI{1e8}{\per\centi\metre\cubed}$, five orders of magnitude below the equilibrium intrinsic concentration $n_i \approx \SI{2.4e13}{\per\centi\metre\cubed}$. Carrier-induced changes to the complex refractive index through free-carrier dispersion are therefore negligible at this power level, validating the separation of the optical and electrical simulations and confirming that the $G_\mathrm{opt}(\mathbf{r})$ computed by FDTD at a single power level can be scaled linearly to any power within the operating range without repeating the electromagnetic computation.

---

## PART IV — SUGGESTED FIGURES AND TABLES

---

**[Suggested Figure — Section 1.1, accompanying Extended Paragraph E1]**

*Description:* A four-panel schematic comparing the square-law detection mechanism for different input field configurations. Panel (a): CW field $E = c$, yielding a constant photocurrent $i = \mathcal{R}c^2$—linear approximation exact. Panel (b): intensity-modulated field $E = \sqrt{P(t)}$ with small modulation depth, showing the photocurrent faithfully tracking $P(t)$. Panel (c): field with separable carrier and signal $E = \tilde{u}(t) + c$ with CSPR $\sim 10$ dB, showing the photocurrent with the SSBI contribution $|\tilde{u}(t)|^2$ visible as a baseband distortion product. Panel (d): block diagram with labels mapping $\tilde{E}(t) \to |\tilde{E}(t)|^2 \to i_\mathrm{ph}(t)$, annotating the three terms of \cref{eq:square_law_expand} and the condition for SSBI negligibility.

*Purpose:* Visualises the domain of validity of the linear responsivity approximation, providing the conceptual foundation for understanding why the high-IQE result of the present chapter is a property of the linear regime and not trivially extensible to high-power non-linear operation.

---

**[Suggested Figure — Section 1.4.3, accompanying Extended Paragraph E2]**

*Description:* Two-panel frequency-response plot. Left panel: normalised photocurrent transfer function $|H_\mathrm{PD}(f)|^2$ in dB plotted from 1 to \SI{300}{\giga\hertz} on a logarithmic frequency axis, showing three curves: $|H_\mathrm{tt}(f)|^2$ (dash-dot sinc$^2$), $|H_\mathrm{RC}(f)|^2$ (dashed single-pole), and their product $|H_\mathrm{PD}(f)|^2$ (solid). The $-3$ dB level is drawn as a horizontal reference, and vertical dashed lines mark $f_\mathrm{tt}$, $f_\mathrm{RC}$, $f_{3\,\mathrm{dB}}$, and $f_\mathrm{Nyq} = R_s/2$ for $R_s = \SI{100}{\giga\baud}$. Right panel: same plot but on a linear frequency scale from 0 to \SI{200}{\giga\hertz}, showing the eye diagram penalty as a shaded region between $f_\mathrm{Nyq}$ and $f_{3\,\mathrm{dB}}$.

*Purpose:* Provides the LTI channel representation of the photodetector that feeds directly into the equaliser design of the DSP chapter. Quantifies the margin between the operating Nyquist frequency and the $-3$ dB roll-off point, confirming the ISI-free regime claimed in Extended Paragraph E2.

---

**[Suggested Figure — Section 1.5, accompanying Extended Paragraph E3]**

*Description:* Two-panel noise contribution plot. Left panel: noise current PSD $i_n^2(f)$ [A$^2$/Hz] in dB versus frequency from 0 to \SI{200}{\giga\hertz} on a linear scale, showing three contributions stacked: $S_\mathrm{shot} = 2qI_d$ (flat blue), $S_\mathrm{thermal} = 4kT/R_L$ (flat green), and a representative $S_\mathrm{TIA}$ with mild $f^2$ roll-up (red). Right panel: total noise variance $\sigma_n^2 = i_n^2 B_e$ as a function of $B_e$ from 1 to \SI{200}{\giga\hertz}, with three curves corresponding to photodetector-only noise (shot + thermal), full assembly noise (adding TIA), and the shot-noise-only floor. Annotated at $B_e = f_{3\,\mathrm{dB}} = \SI{103}{\giga\hertz}$.

*Purpose:* Demonstrates quantitatively that the dark-current shot noise of $I_d = \SI{1.3}{\nano\ampere}$ contributes negligibly to the total $\sigma_n^2$ of the photodetector–TIA assembly at \SI{103}{\giga\hertz}, establishing the correct design prioritisation: TIA noise reduction yields greater sensitivity improvement than further $I_d$ minimisation at this operating point.

---

**[Suggested Figure — Section 4.2, accompanying Extended Paragraph E4]**

*Description:* Three-panel figure. Left panel: required $\mathrm{OMA}_\mathrm{out}$ [dBm] for BER $= 10^{-4}$ as a function of modulation order $M = 2, 4, 8$ at $R_s = \SI{100/\log_2 M}{\giga\baud}$, evaluated from \cref{eq:ser_pam} with $\mathcal{R} = \SI{0.95}{\ampere\per\watt}$ and $i_n = \SI{15}{\pico\ampere\per\sqrt\hertz}$. Each bar shows the breakdown into signal-to-noise and bandwidth contributions. Centre panel: chromatic dispersion transfer function magnitude $|\mathrm{Re}\{H_\mathrm{CD}(f)\}|$ versus electrical frequency $f$ from 0 to \SI{100}{\giga\hertz} for three distances ($z = 1, 10, 100$ km) at O-band ($|\beta_2| = \SI{0.5}{\pico\second\squared\per\kilo\metre}$), showing null-free operation at all intra-datacenter distances. Right panel: the same for C-band ($|\beta_2| = \SI{21}{\pico\second\squared\per\kilo\metre}$) at the same three distances, showing the deep notch falling inside the signal band for $z \geq 1$ km.

*Purpose:* Simultaneously quantifies the PAM-$M$ sensitivity penalty for modulation-order scaling and provides the quantitative dispersion-based justification for O-band operation—connecting device physics (Ge absorption at \SI{1310}{\nano\metre}), modulation format analysis, and system link budget in a single figure. Directly supports the claim in Extended Paragraph E4.

---

**[Suggested Table — Section 4.2, accompanying Extended Paragraph E4]**

*Description:* Four-column table comparing OOK, PAM-4, and PAM-8 for the present device at three operating conditions: (i) shot-noise-limited ($i_n = \sqrt{2qI_d}$), (ii) TIA-limited ($i_n = \SI{15}{\pico\ampere\per\sqrt\hertz}$), and (iii) TIA-limited with $\pm\SI{1}{\decibel}$ responsivity variation. Columns: Modulation format, $R_s$ [Gbaud], $B_e$ [GHz], Required $\mathrm{OMA}_\mathrm{out}$ [dBm], Sensitivity penalty relative to OOK shot-noise limit [dB]. Governing equation column references \cref{eq:ser_pam}.

*Purpose:* Provides a self-contained quantitative bridge between the device-level simulation results of \cref{tab:results} and the system-level sensitivity budget, making the chapter's device-to-system projection directly extensible to multilevel formats without requiring the reader to independently apply \cref{eq:ser_pam}.

---

*End of Chapter Extension Document*

---

**Note on citation integration:** The following reference should be added to the bibliography of the main chapter to support the system-level extensions above:

> D. Pilori, "Digital Signal Processing Techniques for High-Speed Optical Communications Links," Doctoral Dissertation, Politecnico di Torino, 2019. [Cited for: PAM-$M$ BER framework (Eq. (ser\_pam)); chromatic dispersion transfer function derivation; noise current spectral density formulation; equivalent SNR for multi-carrier modulation.]